\documentclass[11pt]{article}

\usepackage{amsmath,amssymb,amsthm}
\usepackage[margin=1in]{geometry}
\usepackage{hyperref}
\newtheorem{lemma}{Lemma}

\title{FO$^k$ Locality: Finite-Type Counting and Corrected Splitting Status}
\author{Inacio F. Vasquez}
\date{}

\begin{document}
\maketitle

\begin{abstract}
We audit finite-variable locality arguments on bounded-degree graphs.
Finite bounded-radius type spaces yield a valid pigeonhole lemma for
\emph{explicitly encoded distinct configurations}, but repetition of
types cannot establish diversity of rooted vertex types. A uniform
local cycle-rank ceiling renders an existential sufficiently-large
threshold formally vacuous, and the vertex-transitive Petersen graph
refutes the intended nonvacuous cycle--orbit splitting implication for
maximum degree three and radius one. A formerly proposed cycle-fan
predicate also needs an explicit ordinary FO$^k$ variable-budget audit;
counting quantifiers must not be silently introduced. The intended
repository-native H1 locality theorem remains open. No general
P-versus-NP, global normalization, or scientific URF closure is claimed.
\end{abstract}

\section{Scope and assumptions}
We work with finite, simple, undirected uncolored graphs unless a
specific statement supplies a different signature. FO$^k$ denotes
ordinary first-order logic with at most $k$ variable names, without
counting, constants, built-in order, or arithmetic. The radius balls
below are induced subgraphs.

\section{Finite configuration-type counting}
\section*{Configuration-Type Pigeonhole: corrected scope}

\noindent\textbf{Status:} The elementary statement below is proved.
A cycle-rank-to-configuration encoding and any consequent splitting of
\emph{different} rooted vertex types remain \textbf{OPEN}.

Fix $k,\Delta,R$ and a finite collection $\mathcal T(k,\Delta,R)$ of
well-defined local configuration types. Finiteness holds when these types
are evaluated on rooted radius-$R$ balls of graphs with bounded degree
and a fixed finite relational vocabulary. No claim about arbitrary
unbounded-radius global FO$^k$ types is intended.

\begin{lemma}[Finite-type pigeonhole]
Let $\mathcal C$ be a finite set of \emph{distinct} configurations with
a specified map $\tau:\mathcal C\to\mathcal T$. If
$|\mathcal C|>|\mathcal T|$, then there exist distinct
$c,c'\in\mathcal C$ with $\tau(c)=\tau(c')$.
\end{lemma}
\begin{proof}
If $\tau$ were injective, the cardinality inequality would be reversed.
\end{proof}

\noindent\textbf{Conditional cycle-basis application.}
If $\dim_{\mathbb F_2}Z_1(B_{2R}(v))>|\mathcal T|$ \emph{and}
there is a proved injection from a basis of that cycle space into
$\mathcal C$ with a specified type map $\tau$, then the lemma gives two
\emph{distinct configurations of equal type}. The injection cannot be
inferred from the dimension inequality alone.

\noindent\textbf{Invalid former inference:} Equal types for two
configurations do not imply \emph{unequal} rooted vertex types at some
pair of vertices. Vertex-transitive graphs, including the Petersen graph,
can have nonzero local cycle rank with only one rooted vertex type.
Therefore the previous ``Consequently'' clause asserting a vertex-type
split is withdrawn; it cannot support a cycle--orbit splitting theorem.


\section{Cycle--orbit splitting obstruction}
\section*{Cycle--Orbit Splitting: corrected status}

\noindent\textbf{Status (November 2026):} The previously asserted
cycle-rank-to-vertex-type-splitting implication is \emph{not proved}.
Its nonvacuous interpretation fails already for degree three and radius one.
The formal existential threshold statement is trivially true if its
threshold exceeds the maximal possible rank; that vacuity is not a locality
or type-splitting theorem.

For a nonempty finite simple graph $G$ of maximum degree $\Delta$, let
$B_{2R}(v)$ denote the induced radius-$2R$ ball and put
\[
\mathsf{CycComp}_R(G)
 = \max_{v\in V(G)}\dim_{\mathbb F_2} Z_1(B_{2R}(v)).
\]
For $\Delta\ge 2$ put
\[
M_{\Delta,R}=1+\Delta\sum_{j=0}^{2R-1}(\Delta-1)^j,
\qquad U_{\Delta,R}=1+\left\lfloor
 \frac{(\Delta-2)M_{\Delta,R}}{2}\right\rfloor .
\]
(The sum is empty when $R=0$.)

\noindent\textbf{Proposition (uniform ceiling).}
For every nonempty graph of maximum degree $\Delta\ge2$,
$\mathsf{CycComp}_R(G)\le U_{\Delta,R}$.
For $\Delta\le1$, every ball is acyclic and the cycle rank is zero.
\begin{proof}
A radius-$2R$ ball is connected and contains at most $M_{\Delta,R}$
vertices by breadth-first counting. With $n$ vertices, it has at most
$\lfloor\Delta n/2\rfloor$ edges. Its cycle-space dimension is therefore
$|E|-n+1\le 1+\lfloor(\Delta-2)n/2\rfloor \le U_{\Delta,R}$.
For $\Delta\le1$ the graph is a disjoint union of vertices and edges.
\end{proof}

\noindent\textbf{Proposition (nonvacuous splitting obstruction).}
Fix $\Delta=3$ and $R=1$. The Petersen graph, represented as the Kneser
graph $KG(5,2)$, has ten vertices, fifteen edges, degree three and diameter
two. Its radius-two ball at any vertex is the whole graph, with
$\mathbb F_2$ cycle rank $15-10+1=6=U_{3,1}$.
Permuting the five underlying symbols acts transitively on its ten
vertices, so \emph{every} isomorphism-invariant rooted vertex type agrees
at every vertex, including every ordinary $\mathrm{FO}^k$ type.
Thus no $t^*(k,3,1)\le6$ forces two different vertex types, whereas
$t^*(k,3,1)>6$ makes the antecedent impossible. This holds for every $k$.

\noindent\textbf{OPEN:} A genuinely nonvacuous type-separation principle
requires independently verified hypotheses that exclude homogeneous
examples, a precise logical signature (ordinary $\mathrm{FO}^k$ versus
counting logic), and an actual refinement operator tied to those semantics.
The former Cycle Fan Predicate / EF Cycle--Linkage Trap reference does
not discharge these obligations. In particular, no type-splitting result
is derived from a pigeonhole repetition statement alone.


\section{Cycle-fan expressibility boundary}
\section*{Cycle Fan Predicate: syntax and semantic boundary}

\noindent\textbf{Status: OPEN for arbitrary fixed $k,R,m$.}
The previously displayed expression containing
$\exists^{\ge m}z$ was described as an ordinary $\mathrm{FO}^k$
formula without proving that it can be expressed using at most $k$
variable names. As written, $\exists^{\ge m}$ is counting-quantifier
notation; counting logic $\mathrm{C}^k$ must not be silently identified
with ordinary $\mathrm{FO}^k$.

For contrast, with at least $m+1$ available variable names the
\emph{simple} property that $x$ has $m$ distinct neighbors is defined in
ordinary first-order logic by
\[
\psi_m(x):=\exists z_1\cdots\exists z_m\,
\Bigl(\bigwedge_{i=1}^{m}E(x,z_i)\ \land\
\bigwedge_{1\le i<j\le m}z_i\ne z_j\Bigr).
\]
This uses $m+1$ names including $x$, and says nothing about independent
cycles or a shared fan. It is \emph{not} a replacement proof for the
cycle-fan claim. Adding bounded paths, fixed anchors, and cycle
independence requires a separate formal construction with its variable
budget and semantic assumptions explicit.

\noindent\textbf{OPEN obligations:} Fix the graph vocabulary and logic;
define an $(m,R)$-fan and independent cycles; construct a corresponding
formula with an audited variable bound; then prove any proposed
cycle-rank-to-fan implication independently. None of those steps follows
from the existing counting-quantifier display.


\section{Discussion and status}
Finite type repetition is proved under its explicit cardinality premise.
The original cycle-rank-to-vertex-type splitting route is not valid as a
nonvacuous general theorem. A code-bound refinement operator, its
semantics, and a correct FO$^k$-admissibility statement remain unproved.
The Petersen obstruction is a classical graph example, not an original
new graph-theoretic result.

\end{document}
